\section[Case study: sandpile dynamics]{Case study: the sandpile Laplacian and the graph of an avalanche}

The Abelian sandpile supplies a compact example in which one declared graph supports several
operators with different state spaces \cite{dhar1990}.  Let a finite graph have a sink and let
$V$ denote its nonsink vertices.  The reduced toppling matrix is
\begin{equation}
\Delta_{ii}=d_i,
\qquad
\Delta_{ij}=-a_{ij}\quad(i\ne j),
\label{eq:book-toppling-matrix}
\end{equation}
where $a_{ij}$ counts edges and $d_i$ includes edges from $i$ to the sink.  A height configuration
$z\in\mathbb Z_{geq0}^{V}$ is stable when $z_i<d_i$ for every $i$.  Toppling an unstable vertex
changes the configuration by
\begin{equation}
T_i z=z-\Delta e_i.
\label{eq:book-single-toppling}
\end{equation}
Repeated toppling ends at a stable configuration because grains can leave through the sink.

Let $\operatorname{Stab}$ denote this stabilization map and define addition followed by
stabilization,
\begin{equation}
a_i(z)=\operatorname{Stab}(z+e_i).
\label{eq:book-sandpile-addition}
\end{equation}
On recurrent configurations the addition operators commute,
$a_ia_j=a_ja_i$, and generate a finite Abelian group.  The number of recurrent configurations is
\begin{equation}
|\mathcal R|=\det\Delta,
\label{eq:book-sandpile-tree-count}
\end{equation}
which also counts spanning trees rooted at the sink by the matrix--tree theorem.

Four graph-like objects can now be audited separately:
\begin{enumerate}
\setlength{\itemsep}{2pt}
\item the substrate graph declares which vertices exchange grains;
\item $\Delta$ is a linear operator on site-height coordinates;
\item $a_i$ is a nonlinear operator on the finite space of recurrent configurations;
\item one realized avalanche has an event-dependent support and causal order.
\end{enumerate}
Equation~\ref{eq:book-sandpile-tree-count} is an exact bridge between the linear operator and the
configuration count.  It does not identify an avalanche with an eigenvector of $\Delta$.  Avalanche
statistics also depend on the driving rule, stabilization, boundary dissipation, and the ensemble
of recurrent states.

This example disciplines the interpretation of protein graph spectra.  A residue-contact graph,
a message-passing Laplacian, and the set of positions affected by one mutation can share vertices
while defining different operators.  For a sequence model, one may define a thresholded mutational
event
\begin{equation}
\mathcal A_\tau(x;j,b)
=\left\{i:
d\!\left(p_i(\cdot\mid x^{j\leftarrow b}),p_i(\cdot\mid x)\right)>\tau
\right\}
\label{eq:book-mutational-avalanche}
\end{equation}
and study its size, geometry, and operator directions.  This is an avalanche-style observable.
A claim of self-organized criticality additionally requires a stationary drive, a dissipation
mechanism, threshold-robust scaling, finite-size analysis across lengths or graph sizes, and nulls
that exclude ordinary heavy-tailed heterogeneity \cite{bak1987}.

The sandpile therefore contributes a constructive warning to spectral analysis: first declare
the state space on which the Laplacian acts, then state which dynamics use it, and finally measure
the event graph generated by that dynamics.  Spectral similarity at the substrate level does not
guarantee similarity of stabilization or avalanche laws.
